\documentclass[10pt,a4paper]{amsart}
\usepackage[margin=19mm]{geometry}
\usepackage{amsmath,amssymb,mathtools}
\usepackage[scr=rsfs,cal=euler]{mathalfa}
\usepackage{xfrac}
\usepackage[unicode,hidelinks,bookmarks=false]{hyperref}
\newcommand{\ie}{i.e.\ }
\newcommand{\cf}{cf.\ }
\newcommand{\resp}{respectively\ }
\newcommand{\Subsection}{\subsection}
\providecommand{\nobreakdash}{\nobreak\hbox{-}}
\setlength{\emergencystretch}{2em}
\multlinegap0pt
\usepackage{bm}
\usepackage{bbm}
\usepackage{dsfont}
\usepackage{xspace}
\usepackage{cancel}
\usepackage{mathtools}
\usepackage{amsthm}
\makeatletter
\@ifundefined{thm}{\newtheorem{thm}{Thm}}{}
\@ifundefined{lem}{\newtheorem{lem}[thm]{Lemma}}{}
\@ifundefined{observation}{\newtheorem{observation}[thm]{Observation}}{}
\makeatother
\newcommand{\meg}{\geqslant}
\newcommand{\mik}{\leqslant}
\title[Forbidden sparse intersections]{Forbidden sparse intersections}
\author{Pandelis Dodos, Miltiadis Karamanlis}
\date{Source: arXiv:2303.16015v4 (CC BY 4.0)}
\begin{document}
\maketitle

\noindent\emph{Source and licence.} Forbidden sparse intersections. Version arXiv:2303.16015v4 (CC BY 4.0), distributed under the Creative Commons Attribution 4.0 International licence, as recorded in the arXiv licence metadata for this exact version and shown on the abstract page https://arxiv.org/abs/2303.16015v4. Full rights notice: licenses/arxiv-2303.16015-CC-BY-4.0.txt.\par\smallskip

\noindent\textit{Editorial scope.} Counted unit: the Lem whose source label is \texttt{l6.1} in the article (source0.tex lines 869--875, source label \texttt{l6.1}), delivered together with the article's own complete original proof of it (source0.tex lines 890--982, one \texttt{proof} environment of 93 lines). No mathematics is written, repaired or re-derived: statement and proof are the article's own text line for line. The proof is detached from the statement in the source: the article states the result at lines 869--875 and prints its proof at lines 890--982, and the proof environment's own header names the statement, so the association is the article's own. Required context retained uncounted: e5.6 (lines 809--811); e6.4 (lines 863--866); o6.2 (lines 878--888). Every definition, display and label of those lines is retained unchanged. Omitted from this package and not redistributed: 1--808, 812--862, 867--868, 876--877, 889--889, 983--1916. Nothing inside a retained excerpt is omitted or altered: every retained body is delivered exactly as the article's lines are written, so no omission receipt of any kind accompanies this package. Citations: the retained excerpts carry no citation command at all, so no bibliography entry of the article is transcribed and no reference is redistributed. Reference adaptation: none was needed and none was made; every reference inside the delivered unit resolves inside it, so no unresolved reference or citation is produced. Printed slips kept literal: no printed word, symbol, hypothesis or number of the counted statement or of its proof was altered. Layout and class: the article's own class is \texttt{amsart} with packages including amssymb, bm, graphicx, amsmath, amsfonts, amsthm, amsbsy, mathtools, mathrsfs, cases, xy, algorithm2e, hyperref, geometry; the wrapper is the lane's pinned \texttt{amsart} editorial wrapper at 10pt on A4, which restates the theorem-like environments the retained text needs and adds the article's own macro definitions that the retained text needs (\texttt{\textbackslash meg}, \texttt{\textbackslash mik}), with extra wrapper packages (bm, bbm, dsfont, xspace, cancel, mathtools). The delivered numbering is the unit's own. Assets: no retained span contains a figure environment, a graphics-inclusion command, a PDF-page inclusion, a table or a TikZ picture; no image file is copied into this package, the package declares no asset dependency and no image file is redistributed with it. Licence: version arXiv:2303.16015v4 of this article is distributed under the Creative Commons Attribution 4.0 International licence, as recorded in the arXiv licence metadata for this exact version and shown on the abstract page https://arxiv.org/abs/2303.16015v4. A full rights notice is delivered at licenses/arxiv-2303.16015-CC-BY-4.0.txt. Characters: every retained excerpt is pure ASCII, so the delivered engine, XeLaTeX with its default Latin Modern text font, reports no missing character.
\begin{equation} \label{e5.6}
S_{d_1}+S_w\mik \ell.
\end{equation}

\begin{equation} \label{e6.4}
1 >(1+\delta)^{S_{d_1}}\, \Big(1+\frac{p}{1-p}\delta\Big)^{S_{d_2}}\,
\Big(1-\delta-2\frac{p}{1-p}\delta^2\Big)^{S_w}\, \exp(-pn\delta^2),
\end{equation}

\begin{observation} \label{o6.2}
The following hold.
\begin{enumerate}
\item[(i)] \label{i} We have $\frac{1}{1-x}=1+\frac{x}{1-x}$ for every $x\neq 1$.
\item[(ii)] \label{ii} We have $x-\frac{x^2}{2}\mik\ln(1+x)\mik x$ for every $x\meg 0$.
\end{enumerate}
In particular, for every $0\mik x <1$ we have
\begin{equation} \label{en.6.1}
\frac{x}{1-x} -\frac{x^2}{2(1-x)^2} \mik \ln\Big(\frac{1}{1-x}\Big) \mik \frac{x}{1-x}.
\end{equation}
\end{observation}

\begin{lem} \label{l6.1}
We have
\begin{align}
\label{e6.5} & \ \ \ \, S_{d_1}-S_w < 5pn\delta, \\
\label{e6.6} & S_{d_2} -\frac{1-p}{p}S_w  < 3n \delta.
\end{align}
\end{lem}

\begin{proof}[Proof of Lemma \emph{\ref{l6.1}}]
We start with the proof of \eqref{e6.5}. Notice first that, by \eqref{e6.4},
\begin{equation} \label{e6.7}
1 > (1+\delta)^{S_{d_1}} \, \Big(1-\delta-2\frac{p}{1-p}\delta^2\Big)^{S_w}\exp(-pn\delta^2)
\end{equation}
or, equivalently,
\begin{equation} \label{e6.8}
1 >(1+\delta)^{S_{d_1}-S_w}
\Big((1+\delta) \big(1-\delta-2\frac{p}{1-p}\delta^2\big)\Big)^{S_w} \exp(-pn\delta^2).
\end{equation}
Since $(1+\delta)(1-\delta-2\frac{p}{1-p}\delta^2)=1-\frac{1+p}{1-p}\delta^2-2\frac{p}{1-p}\delta^3$,
after taking logarithms and rearranging we find that
\begin{equation} \label{e6.9}
S_{d_1}-S_w < \frac{1}{\ln(1+\delta)}
\left(S_w\, \ln\Big(\frac{1}{1-\frac{1+p}{1-p}\delta^2-2\frac{p}{1-p}\delta^3}\Big)+pn\delta^2\right)
\end{equation}
that implies, by Observation \ref{o6.2}, that
\begin{equation} \label{e6.10}
S_{d_1}-S_w < \frac{1}{\delta(1-\frac{\delta}{2})} \left( S_w \, \frac{\frac{1+p}{1-p}\delta^2+2\frac{p}{1-p}\delta^3}{1-\frac{1+p}{1-p}\delta^2-2\frac{p}{1-p}\delta^3}
+pn\delta^2\right).
\end{equation}
By \eqref{e5.6}, the fact that $\ell< pn$ and \eqref{e6.10}, we conclude that
\begin{equation} \label{e6.11}
S_{d_1}-S_w<pn\delta \left(\frac{\frac{1+p}{1-p}+2\frac{p}{1-p}\delta}{\left(1-\frac{\delta}{2}\right)
\left(1-\frac{1+p}{1-p}\delta^2-2\frac{p}{1-p}\delta^3\right)}+\frac{1}{\left(1-\frac{\delta}{2}\right)}\right).
\end{equation}
The desired estimate \eqref{e6.5} follows from \eqref{e6.11} and the fact that $0<\delta<\frac{1}{10}$ and $p\mik \frac12$.

We proceed to show that inequality \eqref{e6.6} is also satisfied.
As before, we first observe that \eqref{e6.4} yields that
\begin{equation} \label{e6.12}
1 > \Big(1+\frac{p}{1-p}\delta\Big)^{S_{d_2}-S_w}
\left(\big(1+\frac{p}{1-p}\delta\big)\big(1-\delta-2\frac{p}{1-p}\delta^2\big)\right)^{S_w}\, \exp(-pn\delta^2).
\end{equation}
On the other hand, since $0<\delta<\frac{1}{10}$, we have
\begin{equation} \label{e6.13}
\Big(1+\frac{p}{1-p}\delta\Big)\, \Big(1-\delta-2\frac{p}{1-p}\delta^2\Big)\meg
1-\frac{1-2p}{1-p}\delta-\frac{16}{5}\frac{p}{1-p}\delta^2
\end{equation}
that combined with \eqref{e6.12} yields that
\begin{equation} \label{e6.14}
1 > \Big(1+\frac{p}{1-p}\delta\Big)^{S_{d_2}-S_w}
\Big(1-\frac{1-2p}{1-p}\delta-\frac{16}{5}\frac{p}{1-p}\delta^2\Big)^{S_w}\, \exp(-pn\delta^2).
\end{equation}
Now after taking logarithms and rearranging, we have
\begin{equation} \label{e6.15}
S_{d_2}-S_w < \frac{1}{\ln\left(1+\frac{p}{1-p}\delta\right)} \,
\left(S_w\, \ln\left(\frac{1}{1-\frac{1-2p}{1-p}\delta-\frac{16}{5}\frac{p}{1-p}\delta^2}\right)+pn\delta^2\right);
\end{equation}
by Observation \ref{o6.2}, this yields that
\begin{equation} \label{e6.16}
S_{d_2}-S_w < \frac{1-p}{p\delta\left(1-\frac{p}{2(1-p)}\delta\right)} \,
\left(S_w\, \frac{\frac{1-2p}{1-p}\delta+\frac{16}{5}
\frac{p}{1-p}\delta^2}{1-\frac{1-2p}{1-p}\delta-\frac{16}{5}\frac{p}{1-p}\delta^2} + pn\delta^2\right)
\end{equation}
that can be further simplified to
\begin{equation} \label{e6.17}
S_{d_2}-S_w <  S_w\, \frac{\frac{1-2p}{p}+\frac{16}{5}\delta}{\left(1-\frac{p}{2(1-p)}\delta\right)
\left(1-\frac{1-2p}{1-p}\delta-\frac{16}{5}\frac{p}{1-p}\delta^2\right)}+
\frac{(1-p)n\delta}{1-\frac{p}{2(1-p)}\delta}.
\end{equation}
On the other hand, since $p\mik \frac12$, for every $0<\delta<\frac{1}{10}$ we have
\begin{equation} \label{e6.18}
\left(1-\frac{p}{2(1-p)}\delta\right)
\left(1-\frac{1-2p}{1-p}\delta-\frac{16}{5}\frac{p}{1-p}\delta^2\right)\meg 1-\delta;
\end{equation}
indeed, after noticing that
\begin{equation} \label{en.6.2}
1-\frac{1-2p}{1-p}\delta-\frac{16}{5}\frac{p}{1-p}\delta^2 =
1-\delta+ \frac{p}{1-p}\delta -\frac{16}{5}\, \frac{p}{1-p}\delta^2,
\end{equation}
the desired estimate \eqref{e6.18} follows from the elementary inequality
\begin{equation} \label{en.6.3}
\frac{p}{1-p}\delta + \frac{16}{5}\, \frac{p^2}{2(1-p)^2}\delta^3
- \frac{p}{2(1-p)}\delta(1-\delta) - \frac{p^2}{2(1-p)^2}\delta^2
- \frac{16}{5}\, \frac{p}{1-p}\delta^2 \meg 0.
\end{equation}
By \eqref{e6.17} and \eqref{e6.18}, we obtain that
\begin{equation} \label{e6.19}
S_{d_2}-S_w < S_w\, \left(\frac{1-2p}{p}+\frac{16}{5}\delta\right)
\left(1+\frac{10}{9}\delta\right)+(1-p)n\delta\left(1+\delta\right).
\end{equation}
We then expand \eqref{e6.19} to
\begin{equation} \label{e6.20}
S_{d_2}-S_w < S_w\, \left(\frac{1-2p}{p}+ \frac{16}{5}\delta+ \frac{1-2p}{p}\frac{10}{9}\delta +\frac{10}{9}\frac{16}{5}\delta^2\right)+(1-p)n\delta\left(1+\delta\right).
\end{equation}
By \eqref{e5.6}, we see that $S_w\mik \ell< pn$, and so \eqref{e6.20} yields that
\begin{equation} \label{e6.21}
S_{d_2}-\frac{1-p}{p}S_w < n\delta \left(p\frac{16}{5}+(1-2p)\frac{10}{9}+p\frac{10}{9}\frac{16}{5}\delta+(1-p)\left(1+\delta\right)\right).
\end{equation}
Inequality \eqref{e6.6} follows from \eqref{e6.21} and the fact that
$0<\delta<\frac{1}{10}$ and $0<p\mik\frac{1}{2}$.
\end{proof}

\end{document}
